\documentclass[conference]{IEEEtran}
\usepackage[utf8]{inputenc}
\usepackage[T1]{fontenc}
\usepackage[english]{babel}
\usepackage{graphicx}
\usepackage{array}
\usepackage{booktabs}
\usepackage{amsmath}
\usepackage{url}
\usepackage{cite}

\begin{document}

\title{\begin{tabular}{@{}m{2.3cm}@{\hspace{0.3cm}}m{\dimexpr\textwidth-3.0cm\relax}@{}}

\includegraphics[width=2.1cm]{figuras/logo_uni.png} &
\centering\arraybackslash Automatic Extraction of Electoral Results
from Vote-Counting Records of the 2026 Peruvian General Elections
\end{tabular}}

\author{
\IEEEauthorblockN{Josemanuel Rossy Ca\~nari Palante, Kenny Asto Hinostroza,
Melissa Dessire Aylas Barranca, Carlos P\'erez P\'erez}
\IEEEauthorblockA{Master's Program in Artificial Intelligence --- Computer Vision, Group 3\\
National University of Engineering, Lima, Peru}
}

\maketitle

\begin{abstract}
In Peru, official electoral results are published through vote-counting records,
many of which are scanned. This hinders the automatic processing of information
at scale. We present a complete system that downloads records from the National
Office of Electoral Processes (ONPE) portal, locates 46 fields in each document
using a template and printed reference marks, and recognizes handwritten counts
with optical character recognition (OCR) restricted to digits. The evaluation
did not require manual labeling: the officially digitized and published ONPE
results were used as the reference. The test set consisted of a random sample
of 100 records from 23 departments. Twenty-four percent of the records were
digital documents with an extractable text layer, achieving 95.1\% field-level
accuracy. For the 76 scanned handwritten records (3,268 fields), the system
achieved 59.6\% field-level accuracy and a Character Error Rate (CER) of
0.401. On the same records, a multimodal Large Language Model (LLM) achieved
48.4\%. The results show that a specialized, multi-stage system is a useful
foundation for improving handwritten-digit recognition in electoral documents.
\end{abstract}

\begin{IEEEkeywords}
document analysis, optical character recognition, handwritten digit recognition,
electoral documents, template alignment
\end{IEEEkeywords}

\section{Introduction}
Vote-counting records contain the official results for every polling station in
Peru. Although ONPE publishes these documents digitally, much of the information
is contained in scanned forms in which polling-station officials wrote the votes
by hand. Reviewing these records manually is slow, can introduce errors, and is
not feasible for the approximately 88,000 polling stations in a general election.
This limits automated auditing and independent verification of the results.

This work addresses the following question: \emph{how accurate can a traditional
computer-vision system be when extracting results from Peruvian records selected
at a national scale, and what does this imply for a real-world system?} To answer
it, we built and evaluated a five-stage process: downloading, image preparation,
field localization, digit recognition, and generation of structured output. The
study used records from the first round of the 2026 General Elections. The main
contributions are as follows:

\begin{enumerate}
\item \textbf{A reproducible evaluation without manual labeling.} The ONPE portal
publishes both the image of each record and the officially digitized counts.
These values were used as the reference for evaluating the system. Thus, each
record provides approximately 43 already-verified field--value pairs. The test
sample was randomly selected nationwide using a published seed.

\item \textbf{A template aligned using printed form marks.} The record has a
fixed layout and contains black squares along its edges. These marks make it
possible to correct displacements and small distortions introduced during
scanning. This correction improved accuracy by 8.7 percentage points compared
with using fixed coordinates without alignment.

\item \textbf{Identification of two record types.} Twenty-four percent of the
sample consists of digitally generated records whose text layer can be directly
extracted and achieves 95.1\% field-level accuracy. The remaining 76\% are
scanned handwritten records that require image processing. Therefore, a
real-world system must first identify the document type and apply the appropriate
process.

\item \textbf{A clear measure of general-purpose OCR performance.} On handwritten
records, digit-restricted OCR achieved 59.6\% field-level accuracy. This result
shows the limitation of using a general recognizer without training it on this
document type and supports the need to train a specialized model.

\item \textbf{A comparison with a multimodal language model.} Under the same
records and evaluation conditions, a multimodal model given the complete page
and two solved examples achieved 48.4\%, compared with 60.9\% for the multi-stage
system. The multimodal model offers a useful initial approach, but it does not
outperform the specialized method for reading handwritten digits.
\end{enumerate}

The code, field-level results, and sampling seed are publicly
available.\footnote{\url{https://github.com/carlosperez100/onpe_actas}}

\section{Related Work}
\textbf{Object detection.} RT-DETR \cite{rtdetr} is a real-time object-detection
model that, according to its authors, outperforms similarly sized YOLO models in
accuracy and speed. In vote-counting records, where cells are close together and
may be covered by stamps or signatures, learned detectors could be an alternative
for locating fields. In this work, field localization is solved with a template
aligned to the document. In the future, this template could also automatically
generate labels to train detectors such as YOLOv11 or RT-DETR.

\textbf{Text recognition.} TrOCR \cite{trocr} uses \emph{transformer} neural
networks for optical character recognition. DTrOCR \cite{dtrocr} proposes a
variant for printed, scene, and handwritten text recognition. These studies
support the use of specialized models to improve handwritten-digit reading. In
our case, it is important to prevent a model from changing a number into a
seemingly more likely sequence, because each value represents an electoral count.

\textbf{End-to-end document analysis.} OmniParser \cite{omniparser} integrates
text localization, key information extraction, and table recognition. Donut
\cite{donut}, by contrast, seeks to generate the document structure directly
from the image. Results from these works show that documents with extensive
tables can be difficult to process in a single step. We therefore retain a
multi-stage process, as it enables clear inspection and measurement of the
performance of each system component. End-to-end methods are considered an
alternative for future work.

The reviewed studies do not specifically evaluate handwritten digits in Peruvian
electoral records containing stamps and signatures. This is the practical need
addressed by this study.

\section{Dataset}
\subsection{Source and download}
All documents come from the official ONPE results portal for the first round of
the 2026 General Elections, held on April 12, 2026. The portal allows users to
query the records for a polling station, their processing status, and associated
files. Each vote-counting record is downloaded as a PDF from a temporary address
provided by the portal. Downloads were performed at a moderate rate and logged
in a file that allows them to be resumed after an interruption
(\texttt{src/scraper/}).

\subsection{Sample design}
The study population contains approximately 88,064 polling-station codes,
estimated through queries to the portal services. A simple random sample without
replacement of 100 records was drawn using a published seed (2026). The sample
covers 23 departments, identified using \emph{ubigeo} prefixes, where \emph{ubigeo}
is the standard geographic code used in Peru. Lima accounted for 36\%, consistent
with its greater weight in the electorate. This design avoids the bias of
selecting polling stations with consecutive codes from a single district.

\subsection{Record types}
Reviewing the number of pages and the PDF text layer identified two document
types: \textbf{76 handwritten records}, which are one-page scans with no
extractable text, and \textbf{24 digital STAE records}, which have two pages, a
digital signature, and counts written in a digital typeface. STAE stands for
\emph{Soluci\'on Tecnol\'ogica de Apoyo al Escrutinio} (Technological Solution
for Supporting Vote Counting), an ONPE tool that helps polling-station officials
record and print the records. In this study, STAE records are the documents
digitally generated using that tool.

The proportion of digital records in the sample is 24\%, with a 95\% Wilson
confidence interval from 16.7\% to 33.2\%. This result confirms that the
population includes both document types, although the sample is not intended to
estimate their exact proportion. Vision experiments were performed using the 76
handwritten records; digital records were processed through direct extraction of
the text layer (Sec.~\ref{sec:stae}).

\subsection{Reference values}
For each polling station in the sample, the counts officially digitized by ONPE
were saved: votes by political organization, blank, null, and challenged votes,
totals, and the number of voters. Positions 04 and 14 of the form correspond to
withdrawn parties; they are therefore empty in the record and have no value in
the official data. They were excluded from the evaluation to avoid counting
matches between empty fields as correct predictions. In total, each record
provides 43 evaluable fields, for a total of 3,268 fields in handwritten records.

\section{Method}
\subsection{Overview}
The system comprises five stages: (1) downloading the PDF and official reference
values; (2) image preparation; (3) region-of-interest localization; (4) value
recognition; and (5) generation of a structured file for each record. Image
preparation includes converting the PDF to an image at 300 pixels per inch,
correcting skew, improving contrast, and reducing scan noise.

\subsection{Field localization using a template and reference marks}
\label{sec:registro}
The record has a fixed format. Therefore, 46 regions were defined once in a
reference record: 38 vote cells, four summary rows, and four header fields. The
coordinates were stored as proportions of page width and height
(Fig.~\ref{fig:overlay}).

The direct use of fixed coordinates fails when the image contains black margins,
displacements, or slight distortions introduced by the scanner. In some records,
the content was shifted, causing a cell to be cropped over the wrong row. To
solve this, the black squares printed along the form edges are used as reference
marks. First, dark squares near the borders are detected. They are then compared
with 15 known positions in the reference record to compute the most consistent
displacement, ignoring missing marks or false detections. Finally, a simple
geometric correction is applied to the page before cropping the fields.

\begin{figure}[t]
\centering
\includegraphics[width=\columnwidth]{figuras/plantilla_overlay.png}
\caption{Reference record with the 46 template regions: 38 vote cells (blue),
summary and header fields (red). Coordinates are applied after aligning the
document using the reference marks.}
\label{fig:overlay}
\end{figure}

\subsection{Recognition of handwritten counts}
Each cell is processed with EasyOCR, configured to recognize digits only. Before
recognition, improvements obtained from an error analysis of 10 pilot records
are applied (Fig.~\ref{fig:limpieza}). First, an inner margin is removed to
prevent the cell's printed borders from being interpreted as digits. The image
is then converted to black and white using a threshold computed for each cell,
so that handwritten ink is highlighted. Dotted separators printed at known cell
positions and small scanning artifacts are also removed. In fields containing
one box per digit, each box is read separately. If this process removes overly
faint strokes, the system retries recognition using the original crop.

\begin{figure}[t]
\centering
\includegraphics[width=\columnwidth]{figuras/limpieza_celdas.png}
\caption{Cleaning of problematic cells (top: original crop; bottom: processed
crop). Borders and dotted separators that could be mistaken for digits are
removed.}
\label{fig:limpieza}
\end{figure}

\subsection{Rule for empty cells}
Polling-station officials often leave blank the cell of a party that received no
votes. In this form, an empty cell represents zero and does not necessarily
indicate an OCR error. Therefore, when no value is recognized in a party cell or
in challenged votes, the system assigns zero. This contribution of form-specific
knowledge is reported separately to distinguish it from the image-recognition
performance itself.

\subsection{Processing of digital records}
\label{sec:stae}
For digital STAE records, the results are extracted directly from the PDF text
layer. Each party name appears next to its count, and summary labels appear
before their values. Extracted names are compared with the official names after
normalizing capitalization, spaces, and accents. Thus, no image processing is
required.

\section{Experiments}
\subsection{Metrics and statistical procedure}
\textbf{Field-level accuracy} is the percentage of fields whose read value
exactly matches the official value. \textbf{CER} (\emph{Character Error Rate})
is the character-level error rate: it compares recognized digit strings with the
correct strings. A lower CER indicates fewer errors.

Fields from the same record may be related because they share the same scan
quality and handwriting. For this reason, confidence intervals were calculated
by resampling complete records, with 10,000 repetitions. Field-level results and
per-record metrics are available in the project repository.

\subsection{Incremental evaluation of improvements in the pilot (10 records)}
In an initial analysis, consecutive polling stations from a single district were
used. Therefore, it does not represent the country, but it does allow us to
measure the contribution of each improvement added to recognition
(Table~\ref{tab:pilot}).

\begin{table}[t]
\caption{Cumulative effect of improvements in the 10-record pilot.}
\label{tab:pilot}
\centering
\begin{tabular}{lrr}
\toprule
Configuration & Accuracy (\%) & CER \\
\midrule
v1: digit OCR without cleaning & 43.95 & 0.524 \\
v2: + ink, border, and box cleaning & 52.09 & 0.419 \\
v3: + separator and artifact removal & 52.79 & 0.476 \\
v4: + empty-cell = zero rule & 58.37 & 0.420 \\
\bottomrule
\end{tabular}
\end{table}

\subsection{National evaluation (76 handwritten records)}
Table~\ref{tab:national} presents the system results on the nationwide random
sample. Alignment through reference marks proved essential. Without it, framing
variations severely affected some records: the worst one achieved 14\% accuracy
and a CER above 1. With alignment, accuracy increased by 8.7 percentage points,
and the worst record rose from 14\% to 30\%. The empty-cell-equals-zero rule
added 21.6 percentage points, as these cases were frequent in the national
sample.

\begin{table}[t]
\caption{Results on the nationwide random sample of 76 handwritten records.}
\label{tab:national}
\centering
\begin{tabular}{lrr}
\toprule
Configuration & Accuracy (\%) & CER \\
\midrule
Fixed template, OCR without cleaning & 26.81 & 0.724 \\
Fixed template + zero rule & 50.86 & 0.482 \\
Aligned template, OCR without cleaning & 38.00 & 0.617 \\
\textbf{Aligned template + zero rule} & \textbf{59.58} & \textbf{0.401} \\
\bottomrule
\end{tabular}
\end{table}

\textbf{Main result.} The system achieved \textbf{59.58\% field-level
accuracy}, with a 95\% confidence interval from 56.61\% to 62.48\%. CER was
\textbf{0.401}, with a 95\% confidence interval from 0.360 to 0.444. Mean
per-record accuracy was 59.6\%, with a standard deviation of 13.0 percentage
points and a range from 30\% to 84\%. The remaining differences are mainly due
to handwriting and scan quality rather than field localization.

\textbf{Digital records.} Text-layer extraction from the 24 STAE records
achieved \textbf{95.12\%} accuracy (955 of 1,004 fields). Remaining errors are
related to matching party names, not to reading numbers. The 35.5 percentage
point difference from handwritten records confirms the importance of identifying
the document type before processing it.

\subsection{Comparison with a multimodal language model}
\label{sec:llm}
As an alternative to the multi-stage process, we evaluated a large language
model with the ability to analyze images and text (LLM). The model received the
complete image of each record and two example records with their structured
outputs. It had to return the counts without using a template, alignment, or
independent OCR. Both methods were compared on the same 15 handwritten records.
The model used was Gemini 2.5 Flash, accessed on its free tier in July 2026.
The code is available in \texttt{src/recognition/llm\_ocr.py}.

\begin{table}[t]
\caption{Comparison between the multi-stage system and a multimodal model with
two examples, on the same 15 handwritten records.}
\label{tab:llm}
\centering
\begin{tabular}{lrr}
\toprule
Method & Accuracy (\%) & CER \\
\midrule
\textbf{Multi-stage system (aligned template + OCR)} & \textbf{60.93} & \textbf{0.395} \\
Multimodal LLM with two examples & 48.37 & 0.591 \\
\bottomrule
\end{tabular}
\end{table}

The multi-stage system outperformed the multimodal model by 12.6 percentage
points. The model correctly identified the form's overall structure, but had
greater difficulty reading handwritten digits in low-quality scans. On the two
records with the worst results, it achieved 20.9\% and 23.3\%, whereas the
multi-stage system achieved 34.9\% and 53.5\%, respectively. This is because
the OCR receives an already isolated cell, whereas the multimodal model must
interpret the entire page and its visual noise. The multimodal model may be
useful for preparing an initial set of labels for a new format, but it does not
replace a specialized recognizer for this task.

\subsection{Error analysis}
Errors were concentrated in four situations: thin digits that OCR does not
detect; handwritten strokes that extend into areas where separators are removed;
confusions between similar shapes, such as 8 and 18 or 7 and 9; and fields made
up of several boxes. These are mainly recognition errors rather than field
localization errors. Therefore, the next step is to train a recognizer specific
to these digits and this form.

\section{Conclusions and Future Work}
The proposed system achieved 59.6\% field-level accuracy on handwritten records
selected nationwide and 95.1\% on digital records. The evaluation was
reproducible and did not require manual labeling, as it used the results
officially digitized by ONPE as the reference.

The study yields four main conclusions. First, before cropping fields with a
template, the document must be aligned with its printed reference marks. Second,
the rule that an empty cell equals zero must be distinguished from the strictly
visual performance of OCR. Third, it is advisable to identify whether the
document is handwritten or digital before selecting the processing method.
Fourth, a general multimodal model requires less initial setup, but achieved
12.6 percentage points less accuracy than the specialized system.

As future work, we propose training a \emph{transformer}-based recognizer with
the field--value pairs obtained from the records (approximately 43 per record
and 43,000 from 1,000 records). We also propose training field detectors, such
as YOLOv11 and RT-DETR \cite{rtdetr}, using labels automatically created with
the aligned template; comparing the multi-stage system with end-to-end document
analysis methods \cite{omniparser,donut}; and using the multimodal model to
support the generation of initial labels that can subsequently be reviewed by a
person.

\begin{thebibliography}{9}
\bibitem{rtdetr} Y.~Zhao, W.~Lv, S.~Xu, J.~Wei, G.~Wang, Q.~Dang, Y.~Liu,
and J.~Chen, ``DETRs beat YOLOs on real-time object detection,'' in
\emph{Proc. IEEE/CVF CVPR}, 2024, pp. 16965--16974.
\bibitem{dtrocr} M.~Fujitake, ``DTrOCR: Decoder-only transformer for optical
character recognition,'' in \emph{Proc. IEEE/CVF WACV}, 2024, pp. 8025--8035.
\bibitem{omniparser} J.~Wan, S.~Song, W.~Yu, Y.~Liu, W.~Cheng, F.~Huang,
X.~Bai, C.~Yao, and Z.~Yang, ``OmniParser: A unified framework for text
spotting, key information extraction and table recognition,'' in \emph{Proc.
IEEE/CVF CVPR}, 2024, pp. 15641--15653.
\bibitem{trocr} M.~Li, T.~Lv, J.~Chen, L.~Cui, Y.~Lu, D.~Florencio, C.~Zhang,
Z.~Li, and F.~Wei, ``TrOCR: Transformer-based optical character recognition
with pre-trained models,'' in \emph{Proc. AAAI}, 2023, pp. 13094--13102.
\bibitem{donut} G.~Kim, T.~Hong, M.~Yim, J.~Nam, J.~Park, J.~Yim, W.~Hwang,
S.~Yun, D.~Han, and S.~Park, ``OCR-free document understanding transformer,''
in \emph{Proc. ECCV}, 2022, pp. 498--517.
\end{thebibliography}

\end{document}